\chapter{Estado del Arte}

\section{Contexto del mercado}

El comercio electrónico en Argentina atraviesa un período de crecimiento y profesionalización sostenida. Según el informe NubeCommerce 2026 de Tiendanube, durante 2025 las tiendas del ecosistema alcanzaron casi 22 millones de órdenes, con un crecimiento del 40,2\% en facturación por encima de la inflación acumulada. La Cámara Argentina de Comercio Electrónico (CACE) reportó una facturación de \$15,3 billones en el primer semestre de 2025, con el canal digital representando el 25\% de las ventas totales del país.

En este contexto, la mayoría de los comercios dispone hoy de datos operativos en volumen creciente, pero la capacidad de transformarlos en decisiones estratégicas sigue siendo escasa, especialmente en el segmento PyME.

\section{Herramientas de analytics descriptivo: el estándar actual}

Las herramientas de análisis más utilizadas por los comercios de e-commerce en Argentina pertenecen al nivel descriptivo. Tiendanube ofrece funcionalidades nativas de estadísticas en sus planes pagos, incluyendo desglose de ventas, comportamiento de productos y métricas de clientes. Sin embargo, esta información se presenta en forma de reportes y gráficos que requieren interpretación por parte del usuario.

Google Analytics, ampliamente integrado en las principales plataformas, ofrece análisis de tráfico y comportamiento, pero no genera recomendaciones comerciales. Shopify Analytics presenta capacidades similares dentro de su ecosistema. Ninguna de estas herramientas indica al comerciante qué acción concreta debe tomar ni en qué momento.

\section{Nubelytics: el competidor local más próximo}

En el mercado argentino, Nubelytics se posiciona como la herramienta más cercana a la propuesta del presente proyecto. La plataforma apunta a transformar datos de e-commerce en acciones concretas, con foco en la identificación de oportunidades en ventas y carritos abandonados.

Sin embargo, su enfoque se centra principalmente en métricas de conversión y recuperación de ventas perdidas, sin abordar de forma integral la gestión de inventario, la estacionalidad por rubro ni un sistema de retroalimentación sobre la efectividad de las recomendaciones ejecutadas.

\section{Herramientas internacionales con analytics prescriptivo}

A nivel internacional existen soluciones que se acercan al nivel prescriptivo, aunque presentan barreras significativas para su adopción en el segmento PyME argentino.

\subsection{Zebra Technologies / Reflexis}

Es una plataforma de analytics prescriptivo orientada al retail que genera recomendaciones específicas del tipo ``aplicá un descuento del 15\% en el producto X durante las próximas dos semanas''. Su enfoque es el más cercano conceptualmente al presente proyecto, pero está diseñada para cadenas de retail de escala empresarial, con costos y complejidad de implementación inaccesibles para PyMEs.

\subsection{Conjura}

Es una plataforma de analytics para e-commerce que conecta datos de ventas, publicidad, clientes e inventario para generar insights accionables. Está orientada a marcas de tamaño medio a grande, opera principalmente en inglés y dentro del ecosistema de plataformas internacionales como Shopify.

\subsection{Polar Analytics}

Unifica datos de múltiples fuentes y genera insights orientados a la acción, con fuerte integración en el ecosistema Shopify. No contempla las particularidades del mercado latinoamericano ni ofrece recomendaciones diferenciadas por rubro.

\subsection{Klaviyo}

Es una herramienta de marketing automation que incorpora recomendaciones personalizadas, pero su foco está exclusivamente en comunicaciones con clientes (email, SMS), sin abordar la gestión operativa del comercio.

\section{Análisis comparativo y brecha identificada}

El relevamiento del mercado permite identificar una brecha concreta y relevante. Las herramientas disponibles en Argentina son descriptivas y no recomiendan acciones. Las herramientas internacionales con capacidades prescriptivas están orientadas a empresas de mayor escala, operan en inglés, no contemplan las particularidades del mercado local y presentan costos incompatibles con el segmento PyME.

Ninguna solución identificada combina los siguientes atributos de forma simultánea:

\begin{enumerate}
    \item Recomendaciones accionables priorizadas.

    \item Adaptación a rubros específicos (gastronomía, indumentaria, electrónica).

    \item Integración con la API de Tiendanube como plataforma dominante en Argentina.

    \item Accesibilidad para comercios de pequeño y mediano tamaño.

    \item Un mecanismo de retroalimentación que permita evaluar y mejorar la calidad de las recomendaciones a lo largo del tiempo.
\end{enumerate}

Esta combinación constituye la propuesta de valor diferencial del presente proyecto y fundamenta su pertinencia tanto académica como práctica.
